\chapter{Modellvergleich, Reflexion und Übungen}

\section{Vergleich der Methoden}

\begin{table}[H]
\centering
\small
\begin{tabular}{p{2.8cm}p{2.4cm}p{3.4cm}p{3.4cm}}
\toprule
\textbf{Methode} & \textbf{Typ} & \textbf{Stärken} & \textbf{Grenzen} \\
\midrule
Entscheidungsbaum & Classification & sehr gut interpretierbar & overfitting-anfällig \\
\midrule
Random Forest & Classification & robust, gute Genauigkeit & weniger gut erklärbar \\
\midrule
SVM & Classification & starke Trennung, gut bei klaren Klassen & Parameterwahl anspruchsvoll \\
\midrule
KNN & Classification & einfach, intuitiv & sensitiv auf Skalierung und \(k\) \\
\midrule
Lineare Regression & Prediction & schnell, gut interpretierbar & nur lineare Zusammenhänge \\
\bottomrule
\end{tabular}
\caption{Überblick über die behandelten Verfahren}
\label{tab:modelle-15}
\end{table}

\section{Notebook-Aufträge (Sek II)}

\begin{myexercise}[Notebook-Parcours]
Bearbeiten Sie die fünf Notebooks in dieser Reihenfolge und dokumentieren Sie
Ihre Ergebnisse strukturiert:
\begin{enumerate}
    \item \href{\notebookbaseurl01_Entscheidungsbaum_Classification.ipynb}{\texttt{01\_Entscheidungsbaum\_Classification.ipynb}}
    \item \href{\notebookbaseurl02_Random_Forest_Classification.ipynb}{\texttt{02\_Random\_Forest\_Classification.ipynb}}
    \item \href{\notebookbaseurl03_SVM_Classification.ipynb}{\texttt{03\_SVM\_Classification.ipynb}}
    \item \href{\notebookbaseurl04_KNN_Classification.ipynb}{\texttt{04\_KNN\_Classification.ipynb}}
    \item \href{\notebookbaseurl05_Regression_Prediction.ipynb}{\texttt{05\_Regression\_Prediction.ipynb}}
\end{enumerate}

\textbf{Abgabe:}
\begin{itemize}
    \item Für jedes Notebook: mindestens eine beantwortete Übungsfrage
    \item Ein kurzer Vergleich: Welches Modell würden Sie für den Schulkontext bevorzugen und warum?
\end{itemize}
\end{myexercise}

\begin{mychallenge}[Transferaufgabe]
Entwerfen Sie ein eigenes Mini-Datenprojekt mit mindestens 30 Zeilen Daten:
\begin{itemize}
    \item Entweder Classification oder Prediction
    \item Datenquelle dokumentieren
    \item Ein Modell trainieren und evaluieren
    \item Ergebnis in einem Plot visualisieren
\end{itemize}
Präsentieren Sie Ihr Projekt in 3 Minuten.
\end{mychallenge}

\section{Lernziele}
\begin{todolist}
    \item Ich kann den Unterschied zwischen \textbf{Classification} und \textbf{Prediction} erklären.
    \item Ich kann die Grundidee von \textbf{Entscheidungsbaum}, \textbf{Random Forest}, \textbf{SVM} und \textbf{KNN} beschreiben.
    \item Ich kann ein einfaches Modell in Python trainieren und eine Vorhersage interpretieren.
    \item Ich kann Modellresultate mit geeigneten Metriken einordnen.
    \item Ich kann Ergebnisse in Diagrammen visualisieren und kritisch reflektieren.
\end{todolist}
